\documentclass[a4paper,12pt,twocolumn]{article}
\usepackage[T1]{fontenc}
\usepackage{xcolor}
\usepackage{lipsum}
\usepackage{graphicx}
\usepackage{array}
\usepackage{booktabs}
\usepackage{amsmath}
\usepackage[parfill]{parskip}
\usepackage[style=apa]{biblatex}
\addbibresource{learnlatex.bib}

\newcommand\kw[1]{\textcolor{blue}{\itshape #1}}

\includeonly{pref, chap1}

\begin{document}
\tableofcontents
\listoffigures
\section{Preface}
The preface text. See \cite{doody}.
\section{Introduction}
The first chapter text.

\section{Title of first section}
Testing hellof jieow jfiopwae jfipoew jiojfiew ojfieopw ijfokps;d jfiokfew f
testing 123 test fsdfsd dsa d fdfdfd ds sd sd s ssssssssss dddddddddd dsds

Something about \kw{apples} and \kw{oranges}.

\begin{center}
	\includegraphics[height=2cm]{example-image}
\end{center}
Text
\begin{center}
	\includegraphics[height = 0.25\textheight]{example-image}
\end{center}
test
\begin{center}
	\includegraphics[width = 1\linewidth]{example-image}
\end{center}

\begin{center}
	\includegraphics[clip, trim = 0 0 50 50]{example-image}
\end{center}

\subsection{Subsection}
yrdyfew fsF e\footnote{testing}. hellof~hel

\subsubsection{subsubsection}
testets

\subsubsection{subsubsection}
testets

\paragraph{subsubsubsection}

\section{Second section}
not italic \emph{italic \emph{not siure} italic}.

\section{Lists}
Lists

Ordered
\begin{enumerate}
	\item \begin{enumerate}
		      \item \begin{enumerate}
			            \item nested
		            \end{enumerate}
	      \end{enumerate}
	\item Second
	\item third
\end{enumerate}

Unordered
\begin{itemize}
	\item \begin{itemize}
		      \item \begin{itemize}
			            \item nested
		            \end{itemize}
	      \end{itemize}
	\item random
	\item stuff
\end{itemize}

testing \textit{test}

\lipsum[1-4]
\begin{figure}[ht]
	\centering
	\includegraphics[width = 1\linewidth]{example-image-a.png}
	\caption{caption}
\end{figure}
\lipsum[5]
\begin{figure}[ht]
	\centering
	\includegraphics[width = 1\linewidth]{example-image-b.png}
	\caption{caption}
\end{figure}
\lipsum[6-10]

\begin{tabular}{lcr}
	Animal & Food  & Size   \\
	dog    & meat  & medium \\
	horse  & hay   & large  \\
	frog   & flies & small  \\
\end{tabular}

\begin{tabular}{cp{5cm}}
	Animal & Description                                                          \\
	dog    & The dog is a member of the genus Canis, which forms part of the
	wolf-like canids, and is the most widely abundant terrestrial
	carnivore.                                                                    \\
	\addlinespace
	cat    & The cat is a domestic species of small carnivorous mammal. It is the
	only domesticated species in the family Felidae and is often referred
	to as the domestic cat to distinguish it from the wild members of the
	family.                                                                       \\
\end{tabular}

\begin{tabular}{*{3}{c}}
	\toprule
	Animal & Food                      & Size   \\
	\midrule
	dog    & meat                      & medium \\
	\cmidrule{1-2}
	horse  & hay                       & large  \\
	\cmidrule(rl){1-1}
	\cmidrule{3-3}
	frog   & flies                     & small  \\
	fuath  & \multicolumn{2}{c}{hello}          \\
	\bottomrule
\end{tabular}

\begin{tabular}{lll}
	\toprule
	\multicolumn{1}{c}{Animal} & \multicolumn{1}{c}{Food}    & \multicolumn{1}{c}{Size} \\
	\midrule
	dog                        & meat                        & medium                   \\
	horse                      & hay                         & large                    \\
	frog                       & flies                       & small                    \\
	fuath                      & \multicolumn{2}{c}{unknown}                            \\
	\bottomrule
\end{tabular}

\subsection{hello}
\label{subsec:labelone}

Text of material for the first subsection.
\begin{equation}
	e^{i\pi}+1 = 0
	\label{eq:labeltwo}
\end{equation}

In subsection~\ref{subsec:labelone} is equation~\ref{eq:labeltwo}
\lipsum[1]


y = mx + c
inline
$y = mx + c$
\[y = mx + c\]
\[a^{b + c} + c_{d}\]
\[\sin(\pi/2) * \log_{10}(100) * \theta^{2} \mu\]
\[y = 2 \sin ^ {2} (\theta)\]
\[
	\int_{-\infty}^{+\infty} e^{-x^2} \, dx
\]
\begin{equation}
	\int_{-\infty}^{+\infty} e^{-x^2} \, dx
\end{equation}
\begin{align*}
	y     & = mx + c               \\
	y     & = mx + c + d           \\
	y     & = mx                   \\
	y + c & = mx \quad\text{hello}
\end{align*}
\lipsum[1]
\[
	\begin{matrix}
		a & b & c \\
		d & e & f
	\end{matrix}
	\quad
	\begin{pmatrix}
		a & b & c \\
		d & e & f
	\end{pmatrix}
	\quad
	\begin{bmatrix}
		a & b & c \\
		d & e & f
	\end{bmatrix}
\]
The matrix $\mathbf{M}$.

\lipsum[3]

{\itshape
	\lipsum[4]
}

The mathematics showcase is from \autocite{Graham1995}, whereas
there is some chemistry in \autocite{Thomas2008}.

Some parenthetical citations: \autocite{Graham1995}
and then \autocite[p.~56]{Thomas2008}.

\autocite{alonso2020deepreinforcementlearningnavigation}

\autocite[See][pp.~45--48]{Graham1995}

Together \autocite{Graham1995,Thomas2008}

\printbibliography

\end{document}
